\scanpage{29}
% archive sheet number: 28
\mspage{\handwritten{12}}
% L1–3
Alors $T(k)$ est dans le centre de $G(k)$ en vertu de BIBLE 6. th.2 , ce qui implique que $T$ est dans le centre de $G$, \deleted{en vertu de Exp VI} \inserted{\handwritten{car $\underline{\mathrm{Centr}}_{G}(T)$ est un sous-schéma fermé de $G$ (Exp \uline{VIII} 6.6.), qui contient les points de $G(k)$, donc est identique à $G$ par le Nullstellensatz, ($G$ étant réduit).}}
\margin{\handwritten{Nullstellensatz}}

% L4–5
Enfin, pour référence ultérieure, signalons la proposition triviale suivante, (dont nous avons fait déjà usage implicitement):

% L6–7
\uline{Proposition} \handwritten{1}.17. \quad Soient $G\handwritten{\supset}H\handwritten{\supset}T$ de\handwritten{s} préschémas en groupes de type fini sur $S$. \deleted{Alors:} \handwritten{Conditions équivalentes:}

% L8–9
\handwritten{(i)} \deleted{\unsure{Si}} $T$ est un tore maximal de $G$, \deleted{$T$ est un tore maximal de $H$, \handwritten{et pour tout $s\in S$, \unsure{$H_s$ et $G_s$ ont même rang réductif|$G_s$ et $H_s$ ont même rang réductif}}}

% L10–12
\handwritten{(ii)} \deleted{\unsure{Si}} $T$ est un tore maximal de $H$, et \deleted{\unsure{si}} pour tout $s\handwritten{\in}S$, on a égalité des rangs réductifs $\handwritten{\rho}_{r}(G_{s})=\handwritten{\rho}_{r}(H_{s})$\handwritten{,} \deleted{alors $T$ est un tore maximal de $G$}.

% L13
\handwritten{2}. \quad \uline{Le groupe de Weyl}.

% L14–19
Soit d'abord $k$ un corps algébriquement clos, et soit $G$ un groupe algébrique sur $k$, \inserted{lisse et affine sur $k$.} Si $T$ est un tore maximal, $C$ son centralisateur et $N$ son normalisateur, alors en vertu de \inserted{\handwritten{Exp \uline{XI}}}5.9. ce sont des sous-groupes fermés \deleted{\unsure{essentiellement}} \inserted{lisses de $G$,} et $\handwritten{C}$ est un sous-groupe ouvert de $N$, de sorte que $\inserted{W=}N/C$ est un groupe fini \deleted{\unsure{étale, qu'on appelle le groupe de Weyl}} \inserted{étale sur $k$ , donc déterminé par le groupe $W(k)$} de \deleted{\unsure{$G$ relatif au tore maximal $T$.}} ses points à valeurs dans $k$, comme $W=W(k)_{k}$ (groupe constant défini par le groupe fini ordinaire $W(k)$).
% L20–25
Le groupe fini $W(k)$ s'appelera le \uline{groupe de Weyl géométrique}, \handwritten{(}ou simplement groupe de Weyl si une confusion n'est pas à craindre\handwritten{)}, \uline{de} $G$ \uline{relativement} \uline{à} $T$. En vertu du théorème de \inserted{conjug\deleted{u}aison \handwritten{de}} Borel, les groupes de Weyl relatifs aux divers tores maximaux sont isomorphes entre eux, c'est pourquoi on parle parfois ``du'' groupe de Weyl de $G$, sans préciser de tore maximal. \deleted{\unsure{En}}

% L26–28
\handwritten{C}omme la formation de $C,N$, et $N/C$ commute à toute extension de la base, on voit que si $k'$ est une extension algébriquement close de $k$, le groupe de Weyl géométrique de $G_{k'}$ relativement à $T_{k'}$\handwritten{\unsure{/}} est canoniquement isomorphe
